\begin{definition}[Realization]
A tuple \(\bar b\in M^n\) \emph{realizes} \(p(\bar x)\) in
\(\mathcal M\) iff \(\mathcal M\entails\varphi(\bar b)\) for every
\(\varphi(\bar x)\in p\).  The structure \(\mathcal M\) \emph{omits}
\(p\) iff no tuple in \(M^n\) realizes \(p\) in \(\mathcal M\).
\end{definition}

\begin{example}[A finitely satisfiable type not realized]
Work in the parameter expansion of \((\N,<)\) having the previously defined
constant \(c_n\) for each \(n\in\N\), and let
\[
  p(x)=\set{c_n<x:n\in\N}.
\]
Thus every displayed member is a formula of the fixed parameter language.
Every finite subset is realized in \(\N\), but
the whole type is not.  Assuming the Axiom of Choice in the ambient set theory,
compactness shows that \(p\) is realized in an elementary
extension of \((\N,<)\).  Such a realization is greater than every standard
natural number.
\end{example}

\begin{definition}[Saturation]
Let \(\kappa\) be an infinite cardinal.  The structure \(\mathcal M\) is
\emph{\(\kappa\)-saturated} iff, whenever \(A\subseteq M\) has
\(\abs{A}<\kappa\), every complete type in finitely many variables over
\(A\) consistent with \(\mathcal M\) is realized in \(\mathcal M\).
An infinite structure is \emph{saturated} iff it is
\(\abs{M}\)-saturated.  This reference does not apply the unqualified term
``saturated'' to a finite structure, because \(\kappa\)-saturation has just
been defined only for infinite \(\kappa\).
\end{definition}

\begin{remark}
Some sources allow tuples of length less than \(\kappa\) in the definition of
\(\kappa\)-saturation.  This reference uses finite tuples.  The convention must
be restated before applying a theorem quoted from a source using the stronger
version.
\end{remark}

\begin{theorem}[Existence of saturated elementary extensions]
\label{thm:saturated-extension}
Assume the Axiom of Choice in the ambient set theory.  Let \(\mathcal M\) be a
structure and let \(\kappa\) be an infinite cardinal.
There is a \(\kappa\)-saturated elementary extension
\(\mathcal N\succcurlyeq\mathcal M\).
\textit{Status: cited theorem.}
\end{theorem}

See \cite{Marker2002} for a complete proof.
The construction iterates elementary extensions realizing all types over
small parameter sets, and takes unions at limit stages using
\cref{thm:elementary-chain}.  Producing a specified cardinal bound requires
additional cardinal-arithmetic bookkeeping and is not asserted here.

